\section{Solution Concepts}\label{sec:solutions}

\subsection{Fixed methods}

With fixed methods \pCBa{}, the planner first decomposes $\tnI$ into a primitive network and then
executes it. A primitive network $\tn$ is \emph{achievable} if it arises from $\tnI$ by finitely
many decompositions. For an achievable $\tn$ there are four ways of executing it.

\begin{itemize}
  \item \emph{Weak LD}: some linearisation of $\tn$ is executable for some choice of outcomes.
  \item \emph{Strong LD}: some linearisation $t_1\cdots t_n$ is executable for every choice of
    outcomes, i.e.\ every $t_i$ is applicable in every state that the prefix $t_1\cdots t_{i-1}$
    can lead to. This is a conformant plan.
  \item \emph{Weak OD}: some policy, choosing the next task depending on the tasks executed so far
    and the current state, executes all of $\tn$ for some choice of outcomes.
  \item \emph{Strong OD}: some such policy executes all of $\tn$ for every choice of outcomes.
\end{itemize}

A problem has a solution of one of these kinds if some achievable primitive network does.

Outcome-dependent execution is best described recursively. Call a pair $(\tn,s)$ of a primitive
network and a state \emph{(strongly) solvable} if $\tn$ is empty, or some unconstrained task
$t$ of $\tn$ is applicable in $s$ and $(\tn\setminus t,s')$ is solvable for every outcome
$s'\in\Out(\alpha(t),s)$. Replacing ``every'' by ``some'' gives \emph{weak} solvability.
These are least fixpoints: solvability is witnessed by a finite tree.

\begin{proposition}\label{prop:od-belief}
  A primitive network $\tn$ has a strong (weak) OD policy from $s$ if and only if $(\tn,s)$ is
  strongly (weakly) solvable.
\end{proposition}

\begin{proof}
  A policy is a function from configurations $(T',s)$, consisting of the set of executed tasks
  and the current state, to the next task. Given a solvable pair, let the policy pick at
  $(T',s)$ a task that witnesses solvability of the remaining network $\tn\setminus T'$ in $s$.
  Every run of this policy then stays among solvable pairs and ends with the empty network.
  Conversely, by induction on the number of remaining tasks, if every run of a policy from
  $(T',s)$ completes, then $(\tn\setminus T',s)$ is solvable, witnessed by the task the policy
  chooses.
\end{proof}

Strong OD solvability in this form is exactly the solvability of \emph{belief states} used by
\DB{}, so their notion and that of \CBa{} agree. We note that the definition of \CBa{}, which
goes through an execution structure and its execution sequences, cannot be taken literally. The
structure forgets along which history a state was reached, and every prefix of an execution
sequence is again an execution sequence. So the requirement that ``every execution sequence is
complete'' fails for every non-trivial network. The recursive formulation is what the proofs of
\CBa{} use, and it is the one we adopt.

\begin{proposition}[\CBa, Props.~2.12, 2.13, 2.15]\label{prop:fm-relations}
  \begin{enumerate}[label=(\roman*)]
    \item A problem has a weak LD solution iff it has a weak OD solution.
    \item Every strong LD solution is a strong OD solution, but not conversely.
    \item For totally ordered problems, strong LD and strong OD solutions coincide.
  \end{enumerate}
\end{proposition}

\begin{proof}
  (i)~A linearisation together with a choice of outcomes is a single run, and a run is what a
  weak policy needs. Formally, both notions are equivalent to weak solvability: a linearisation
  of $\tn$ is an unconstrained task followed by a linearisation of the rest.

  (ii)~Let $S_i$ be the set of states that the prefix $t_1\cdots t_{i-1}$ of a strong LD
  linearisation can lead to. Every state reachable by the policy that follows the linearisation
  lies in the corresponding $S_i$, where the next task is applicable; this gives solvability by
  induction from the end. For the converse, let $a$ add fact $1$ or fact $2$, let $b$ require $1$
  and add $2$, let $c$ require $2$ and add $1$, and let $a\prec b$ and $a\prec c$. The policy that
  runs $b$ first if $1$ holds and $c$ first otherwise succeeds, but each of the two
  linearisations $a\,b\,c$ and $a\,c\,b$ fails for one outcome of $a$.

  (iii)~Decomposition preserves total order, so every achievable network is a chain, and a chain
  has exactly one unconstrained task at every point. A solvable pair therefore forces the same
  task in all states, which is a conformant linearisation.
\end{proof}

\subsection{Method-based policies}

Method-based policies \pCBb{} interleave decomposition and execution. A \emph{node} is a pair
$(\tn,s)$ of a (not necessarily primitive) network and a state. A \emph{policy} maps nodes to
instructions: either an unconstrained primitive task $t$, or a compound task $t$ together with a
method for $\alpha(t)$. It must give the same instruction, up to isomorphism, at isomorphic nodes.
Executing a primitive task leads to $(\tn\setminus t,s')$ for each outcome $s'$; decomposing
leads to the decomposed network in the same state. The nodes reachable from $(\tnI,\sI)$, taken
up to isomorphism, form the \emph{execution structure} of the policy. A node is a \emph{goal} if
its network is empty, and \emph{terminal} if the policy gives no applicable instruction.

\begin{definition}[\CBb, Def.~7]\label{def:mp}
  A policy is
  \begin{itemize}
    \item \emph{weak} if its execution structure is finite and some terminal node is a goal;
    \item \emph{strong} if its execution structure is finite, acyclic, and every terminal node is
      a goal;
    \item \emph{strong cyclic} if every terminal node is a goal.
  \end{itemize}
\end{definition}

We have made two harmless details explicit: executed primitive tasks must be unconstrained, as
in all algorithms of \CBb{}, and an inapplicable instruction makes a node terminal.

\subsection{Strong cyclic solutions and fairness}\label{sec:sc}

The definition of strong cyclic policies contains no condition that would make the execution
reach a goal. Consider a deterministic problem with one compound task $c$, whose only method is
$c\to c$. No primitive network can be reached, so the problem has no solution of any other kind.
But the policy that always decomposes $c$ has an execution structure with a single node and no
terminal node, so it is strong cyclic in the sense of Definition~\ref{def:mp}. More generally,
under this literal reading a deterministic recursive problem may be ``solved'' by refining
forever. Strong cyclic plan existence then does not coincide with plan existence on deterministic
problems, and the deterministic lower bounds that \CBb{} transfer to strong cyclic solutions do
not hold.

The intended notion is clearly that of \textcite{cimatti:et:al:aij-03}: from every node of the
execution structure some goal node must be reachable. If the environment is fair, i.e.\ each
outcome of an action that is executed infinitely often in the same node is eventually taken, such
a policy reaches a goal with certainty.

\begin{definition}\label{def:fair}
  A policy is \emph{(fair) strong cyclic} if from every node of its execution structure some goal
  node is reachable within the execution structure.
\end{definition}

Throughout the rest of this article, ``strong cyclic'' refers to Definition~\ref{def:fair}. It
has the expected properties.

\begin{proposition}\label{prop:fair-basic}
  \begin{enumerate}[label=(\roman*)]
    \item Every strong policy is fair strong cyclic.
    \item For acyclic problems, a fair strong cyclic solution exists iff a strong solution does.
    \item For deterministic problems, a fair strong cyclic solution exists iff a plan exists.
  \end{enumerate}
\end{proposition}

\begin{proof}
  (i)~In a finite acyclic structure every maximal path ends in a terminal node, which is a goal.
  (ii)~In an acyclic problem every refinement terminates, so the execution structure of any
  policy has no infinite paths and no cycles (decomposition makes a network strictly smaller in
  a well-founded order). A fair policy is then strong.
  (iii)~If a plan exists, the policy that follows it is strong. Conversely, the path from the
  initial node to a goal node is a plan.
\end{proof}

\subsection{Games}\label{sec:games}

Both kinds of policies have convenient game characterisations, which underlie all our membership
proofs. They rest on the following observation, which the algorithms of \CBb{} use without
stating it.

\begin{lemma}[delaying decompositions]\label{lem:delay}
  If a problem has a strong (fair strong cyclic) method-based solution, it has one that
  decomposes only unconstrained compound tasks.
\end{lemma}

\begin{proof}
  A decomposition of a task $t$ that is still preceded by other tasks commutes with the
  execution and decomposition of those tasks. Postponing it until $t$ becomes unconstrained gives
  a policy whose runs are those of the original policy with some steps moved later. The later
  decision may depend on more observations, which only helps. For fair policies one checks in
  addition that reachability of a goal is preserved, which follows because every run of the new
  policy can be matched by a run of the old one with the same outcomes.
\end{proof}

With progression only, a strong solution is the same as a winning strategy in an AND/OR game.

\begin{theorem}\label{thm:mp-andor}
  A problem has a strong method-based solution iff the planner wins the following game from
  $(\tnI,\sI)$: a node with the empty network is won; at a node $(\tn,s)$ the planner either picks
  an unconstrained primitive task applicable in $s$, after which the environment picks an
  outcome, or picks an unconstrained compound task and a method to decompose it.
  The planner must win in finitely many moves.
\end{theorem}

\begin{proof}
  A strong policy has a finite acyclic execution structure, so every run is shorter than the
  number of its nodes, and the policy is a winning strategy. Conversely, rank each won node by the
  least number of moves in which it can be won, which is invariant under isomorphism. The policy
  that always moves to a node of smaller rank is well defined up to isomorphism, every run strictly
  decreases the rank, and every node has finitely many successors up to isomorphism. Hence its
  execution structure is finite and acyclic.
\end{proof}

For fair strong cyclic policies the game has a nested fixpoint. A set $X$ of nodes is
\emph{fair-closed} if every node of $X$ is in the \emph{attractor} of $X$. The attractor is the
least set $Z$ containing the goal nodes and such that a node is in $Z$ if
\begin{itemize}
  \item the planner can execute an applicable primitive task all of whose outcomes are in $X$ and
    at least one of which is in $Z$, or
  \item the planner can decompose a task such that the result is in $X$ and in $Z$.
\end{itemize}

\begin{theorem}\label{thm:mp-fair-game}
  A problem has a fair strong cyclic solution iff $(\tnI,\sI)$ lies in a fair-closed set of nodes
  that is closed under isomorphism.
\end{theorem}

\begin{proof}
  The nodes of the execution structure of a fair policy form such a set: the policy's
  instruction keeps all successors in the structure, and a shortest path to a goal shows that each
  node is in the attractor. Conversely, given a fair-closed set, let the policy choose at each
  node a move that witnesses attractor membership of minimal rank. Then all successors stay in
  the set, and following the witnessing outcomes the rank decreases until a goal is reached.
\end{proof}

\subsection{Comparing the two families}

The following facts relate the two families and are used repeatedly.

\begin{proposition}[\CBb, Props.~1, 2 and Fig.~1]\label{prop:fm-vs-mp}
  \begin{enumerate}[label=(\roman*)]
    \item For primitive problems, strong and fair strong cyclic method-based solutions coincide
      with strong OD solutions.
    \item A problem has a weak method-based solution iff it has a weak fixed-method solution.
    \item There are totally ordered problems with a strong method-based solution but no strong
      fixed-method solution.
  \end{enumerate}
\end{proposition}

\begin{proof}
  (i)~Without compound tasks the game of Theorem~\ref{thm:mp-andor} is the recursive definition of
  strong solvability, and every run is finite.
  (ii)~A weak policy has a run ending in a goal; its decompositions give an achievable network
  and its executed tasks a linearisation. Conversely, a weak plan can be turned into a policy by
  following a shortest run of it. A shortest run visits no two isomorphic nodes, so the policy is
  well defined.
  (iii)~Let $\tnI=\langle a,c\rangle$, where $a$ adds fact $1$ or fact $2$, and $c$ has the two
  methods $c\to b$ and $c\to b'$, with $b$ requiring $1$ and $b'$ requiring $2$. A policy can
  decompose $c$ after observing the outcome of $a$, but a fixed decomposition must choose $b$ or
  $b'$ in advance.
\end{proof}
